% Final
\section{Postup při realizaci projektu}

% Final
V této části textu nejprve popíšu úvodní přípravu projektu a jednotlivé kroky, které jsem provedl před zahájením implementace. Následně popíšu, jakým způsobem jsem postupoval při samotné implementaci a jak jsem spolupracoval s kolegou Bc. Jakubem Vondráčkem, který se zabývá backendovou částí této aplikace.

% Final
\subsection{Příprava implementace}

% Final
Před zahájením samotné implementace bylo třeba provést úvodní přípravu projektu. Tuto přípravu bylo třeba provést hned z několika důvodů. Původní aplikace byla naposledy aktivně vyvíjena v roce 2024, kdy jsem ji společně s kolegou Bc. Janem Hlaváčem vyvíjel v rámci našich bakalářských prací. Přestože jsme se projektu věnovali i po dokončení bakalářských prací, naše zaměření bylo především na části projektu které se netýkaly samotného vývoje. Vývoj tak ustál a aplikace nedostávala žádné nové aktualizace. Od roku 2024 došlo k řadě změn, kterým je nyní třeba projekt přizpůsobit.

% Final
První změnou jsou nové verze používaných knihoven. Většina používaných knihoven od posledního vývoje vydala nové verze, které obsahují podstatné opravy a vylepšení. Jejich aktualizace tak mohou přinést řadu výhod, jako například nové funkce, efektivnější vývoj, větší bezpečnost a stabilitu funkcí.

% Final
Další zásadní změnou je nástup nástrojů umělé inteligence v oblasti softwarového inženýrství. Tyto nástroje zásadně mění, jakým způsobem lze postupovat při vývoji softwaru a umožňují vyvíjet aplikace efektivnějším způsobem. Proto je třeba projekt přizpůsobit tak, aby byl připraven pro práci s těmito nástroji.

% Final
Nakonec, došlo k zásadní změně domény, na kterou se tato vzdělávací aplikace zaměřuje. Původní aplikace byla zaměřena výhradně na výuku jazyků. Nová aplikace by měla fungovat jako obecnější rozšiřitelná vzdělávací aplikace se zaměřením na výuku programování. Z tohoto důvodu je třeba změnit strukturu některých částí aplikace tak, aby bylo možné ji v budoucnu lépe rozšiřovat a udržovat. Se změnou domény se také pojí odstranění některých zastaralých funkcí, které neodpovídají novému zaměření aplikace.

% Final
V následujících částech textu tak podrobněji popíšu jednotlivé kroky přípravy projektu. Tuto přípravu jsem prováděl současně s předchozími fázemi analýzy a návrhu.

% Final
\subsubsection{Příprava knihoven komponentů}

% Final
V rámci přípravy aplikace jsem se rozhodl od aplikace oddělit některé UI komponenty a přesunout je do samostatných knihoven, které bude tato aplikace využívat. Komponenty jsem rozdělil do dvou knihoven \texttt{@eduriam/ui-core} a \texttt{@eduriam/ui-x}. První zmíněná knihovna bude obsahovat základní UI komponenty, jako tlačítka, textová pole, karty apod. Druhá zmíněná knihovna pak bude obsahovat komplexnější komponenty, které budou využívány jen ve specifických případech. Toto rozdělení tak reflektuje strukturu knihovny Material UI \cite{MaterialUI}, která je podobně rozdělena na knihovnu základních komponentů MUI Core a na knihovnu komplexnějších komponentů MUI X.

% Final
Důvodů pro rozdělení aplikace do knihoven je hned několik. Vytvoření knihovny \texttt{@eduriam/ui-core} umožní postupně vytvořit jednotný design systém a knihovnu základních komponentů, kterou bude možné využívat jak v samotné aplikaci, tak v dalších částech projektu, jako například na webových stránkách, nebo v dalších interních programech. Díky tomu bude možné dodržet konzistentní design napříč všemi částmi projektu. Dalším důvodem je dodržení design principu \textit{Separation of Concerns} \cite[36]{ADPArchitectures}, který říká, že bychom při vytváření systémů měli jasně oddělit zodpovědnosti jednotlivých částí systému. Díky tomu bude možné docílit přehlednější struktury kódu a lepší testovatelnosti a udržitelnosti celého systému. V případě této aplikace jsem se tak snažil oddělit základní UI komponenty od logiky samotné aplikace.

% Final
Cílem zmíněné knihovny tak bylo oddělit základní UI komponenty od zbytku projektu. Kromě základních komponentů ovšem vznikla potřeba sdílet určité složitější komponenty napříč různými částmi projektu. Zejména se jedná o komponenty, které vizualizují průběh učení dané lekce, tedy komponenty zobrazující jednotlivá cvičení a vysvětlení látky. Tyto komponenty budou využívány jednak v samotné aplikaci a jednak v interním nástroji pro tvorbu obsahu, ve kterém bude možné například zobrazit náhled učení dané lekce. Zmíněný nástroj již není součástí této práce. Jelikož bude knihovna \texttt{@eduriam/ui-core} se základními komponenty využívána ve více programech, snažil jsem se tuto knihovnu zachovat co nejjednodušší a nejmenší. Z tohoto důvodu jsem pro komplexnější komponenty vytvořil samostatnou knihovnu \texttt{@eduriam/ui-x}, kterou bude možné využít v programech, ve kterých budou tyto specifické komponenty skutečně potřeba.

% Final
Další výhodou tohoto přístupu je možnost verzování knihoven pomocí sémantického verzování \cite{SemVer}. Díky tomuto způsobu verzování je možné snadno rozšiřovat a vylepšovat komponenty nezávisle na ostatních částech systému. Knihovny je pak možné postupně aktualizovat a zavádět jejich nové verze do jednotlivých částí projektu. Celý vývoj komponentů je tak strukturovanější a přehlednější.

% Final
Frontend aplikace tak bude využívat obě zmíněné knihovny. Do obou knihoven jsem také zavedl nástroj Storybook\footnote{Dostupné na: \href{https://storybook.js.org/}{https://storybook.js.org/}}, který umožňuje snadno vyvíjet, dokumentovat a testovat UI komponenty v izolovaném prostředí.

% Final
\subsubsection{Aktualizace knihoven a nástrojů}

% Final
Dalším krokem přípravy aplikace byla aktualizace používaných knihoven. Jak jsem již zmínil, v průběhu let byla vydána řada nových verzí používaných knihoven, které mohou přinést řadu dříve zmíněných výhod. Proto jsem se rozhodl aktualizovat nejpodstatnější knihovny a nástroje, které aplikace využívá.

% Final
Mezi aktualizované knihovny patří například Next.js, Material UI, Storybook, ESLint, Prettier, Yarn, Node.js a další. Kromě těchto knihoven jsem dále v průběhu vývoje aktualizoval knihovny, u kterých bylo třeba přejít na jejich novější verze.

% Final
Celkově aktualizace knihoven skutečně přinesla několik výhod. Nejnovější verze knihoven zefektivnily vývoj aplikace, což bylo možné pozorovat například u nástroje Storybook, jehož novější verze přinesla přehlednější a výrazně rychlejší tvorbu UI komponentů díky novým konvencím a rychlejšímu sestavování jednotlivých komponentů při spouštění Storybook serveru.

% Final
\subsubsection{Refaktoring kódu}

% Final
Po aktualizaci knihoven jsem dále provedl refaktoring stávajícího kódu. Nejprve jsem provedl úvodní refaktoring před fází implementace, který se týkal zejména přesunu příslušných komponentů do dříve zmíněných knihoven \texttt{@eduriam/ui-core} a \texttt{@eduriam/ui-x}. Úvodní refaktoring dále zahrnoval zásadní změnu struktury komponentů pro učení, v rámci které jsem komponenty změnil tak, aby byly rozšiřitelnější a lépe udržovatelné. Komponenty jsem tak připravil pro budoucí úpravy a vývoj.

% Final
V rámci úvodního refaktoringu jsem dále vytvořil několik nových komponentů, které usnadňují vývoj a sjednocují rozložení prvků UI na jednotlivých stránkách aplikace. Mezi tyto komponenty patří například \texttt{PageRoot}, \texttt{ContentContainer} a \texttt{PageNavigation}, díky kterým jsou rozložení prvků a navigace napříč stránkami aplikace konzistentní a snadněji konfigurovatelné.

% Final
V průběhu vývoje jsem pak průběžně prováděl další refaktoring, který se týkal zejména odstraňování a aktualizace zastaralých komponentů a stránek. Tyto změny jsem prováděl průběžně zejména kvůli tomu, že bylo snazší staré funkce postupně nahrazovat novými a aktualizovat je, než všechny odstranit najednou.

% Final
\subsubsection{Aktualizace linterů a konvencí}

% Final
V rámci úvodní přípravy projektu jsem také provedl aktualizaci konfigurací nástrojů pro formátování a kontrolu kvality kódu ESLint a Prettier. Jak jsem již zmínil výše, tyto nástroje jsem aktualizoval na jejich novější verze a dále jsem zavedl přísnější pravidla pro kontrolu kódu. Jedním z nově zavedených pravidel je například pravidlo \texttt{eqeqeq}\footnote{Dokumentace dostupná na: \href{https://eslint.org/docs/latest/rules/eqeqeq}{https://eslint.org/docs/latest/rules/eqeqeq}}, které vynucuje používání typově bezpečného porovnávání hodnot pomocí \texttt{===} a \texttt{!==} namísto porovnávání pomocí \texttt{==} a \texttt{!=}.

% Final
Zmíněná pravidla jsem zavedl s cílem snížit riziko vzniku nejednoznačného kódu a chyb, které mohou nastávat při nesprávném použití jazyků JavaScript a TypeScript, ve kterých je frontendová část aplikace naprogramována.

% Final
\subsubsection{Zavedení nástrojů umělé inteligence}

% Final
Další částí přípravy bylo zavedení konvencí pro nástroje umělé inteligence, které mohou vývojářům pomoci s vytvářením kódu. Vycházel jsem ze standardu Agent Skills\footnote{Dostupné na: \href{https://agentskills.io/}{https://agentskills.io/}} definovaným společností Anthropic. Tento standard umožňuje s pomocí souborů ve formátu Markdown definovat instrukce a pravidla pro AI agenty, kteří na základě nich mohou v kódu vykonávat různé úkony. \cite{AgentSkills}

% Final
Zadefinoval jsem tak pravidla a instrukce pro tvorbu nových UI komponentů, nových stránek aplikace, tvorbu automatizovaných testů a práci s \texttt{.feature} soubory, ve kterých jsou definované funkce aplikace.

% Final
Výhodou využití zmíněného standardu je vysoká kompatibilita s nejpoužívanějšími nástroji umělé inteligence pro tvorbu kódu, kterými jsou například Claude Code, Cursor, Codex, Gemini CLI a další. Instrukce tak bude možné využívat napříč různými IDE, nástroji a AI modely. \cite{AgentSkills}

% Final
\subsubsection{Příprava nového návrhu v nástroji Figma}

% Final
V rámci úvodních příprav projektu, které probíhaly současně s analýzou a návrhem, jsem se rozhodl změnit způsob, kterým je tvořen návrh uživatelského rozhraní v nástroji Figma. Návrh uživatelského rozhraní původní aplikace, který jsem vytvořil v rámci mé bakalářské práce, využíval externí knihovnu základních komponentů, které jsem upravil a následně využil v návrhu jednotlivých stránek aplikace.

% Final
Využití knihovny s již připravenými komponenty zpočátku umožnilo rychle a efektivně vytvořit základní prototyp uživatelského rozhraní. Časem se však ukázalo, že má tento přístup zásadní nevýhody při tvorbě \textit{high fidelity} prototypu. Komponenty externí knihovny často využívaly jinou strukturu a parametry, než bylo při návrhu aplikace potřeba. Struktura komponentů a jejich parametrů se tak postupnými úpravami stávala zbytečně složitou, což značně snižovalo efektivitu tvorby prototypu.

% Final
Jelikož mám v rámci této práce za cíl vytvořit kompletně nový a propracovanější vzhled všech komponentů a stránek aplikace, rozhodl jsem se zmíněný přístup změnit a místo externí knihovny komponentů vytvořit všechny komponenty ručně. Díky tomu jsem získal vyšší kontrolu nad tvorbou jednotlivých komponentů a byl jsem schopen je navrhovat jednodušeji a přesněji. Důsledkem toho byl přesnější návrh a snížená složitost komponentů. To vedlo k rychlejším úpravám prototypu a následně k efektivnější implementaci.

% Final
V rámci těchto změn jsem také provedl organizaci návrhů stránek a komponentů, aby odpovídaly nové struktuře aplikace a dříve zmíněným knihovnám \texttt{@eduriam/ui-core} a \texttt{@eduriam/ui-x}.

% Final
Díky výše popsaným změnám by tak měl být návrh uživatelského rozhraní a prototypu v nástroji Figma lépe rozšiřitelný a upravitelný.
